% ARCHIVED v3.84: the corridor block at v3.83 (S&C columns, steps-time frontier before/after with the near bar).
\subsection{UAV-corridor}\label{sec:res:uav:corridor}

The two constraints are read one at a time, as the avoiding geometries are: every table of this scene has a
\emph{tilt} block and a \emph{hump} block, and nothing is averaged across them. A flight succeeds when it
reaches the end of the corridor, the scene's finish line (\autoref{sec:setup:metrics:uav}), and every rate
is the mean over the ten flights of a cell with the sample standard deviation over them, as on the
single-seed tables of D3IL-avoiding.

\subsubsection{Generative Models}\label{sec:res:uav:corridor:raw}
Without projection every model reaches the end of the corridor on every flight at every budget, and none
does so without violating the constraint (\autoref{tab:uav-corridor-raw}): a flight violates the tilt on
$74.5$ to $99.4$ control steps on average and the hump on $28.6$ to $36.6$. The plans follow the level
lanes of the demonstrations, which were flown before either boundary existed
(\autoref{sec:setup:data:uav}), so the two blocks are the same flights read against two constraints: they
share their steps, their time per action differs by at most $1.5$\,ms between the two runs, and they differ
in what they violate. The flow-based models lie within
one standard deviation of one another in violating steps; the diffusion baseline has the most under the
tilt and the fewest over the hump. The flights take $242$ to $251$ control steps to the end of the
corridor, and one action costs $9.1$ to $9.4$\,ms at one evaluation, $170.8$ to $172.3$\,ms at twenty and
about $178$\,ms for the baseline.

\begin{table}[htpb]
  \caption[UAV-corridor: the generative models before projection]{UAV-corridor with the test-time
  constraints, \textbf{before projection}: the plan flown as generated, ten flights per row spread over the
  three routes, one training seed, one block per constraint. The grid: the two average-velocity models at
  $\nfe\in\{1,2,3\}$, instantaneous-velocity matching at $\nfe\in\{1,2,3,5,20\}$, the diffusion baseline at
  its trained budget. Columns as in \autoref{tab:avoiding-raw-models}, each mean $\pm$ the sample standard
  deviation over the flights; \emph{Success}: reaches the end of the corridor
  (\autoref{sec:setup:metrics:uav}); \emph{Steps}: control steps to it. \baselinemark\ the diffusion model
  of \ac{DPCC}.}
  \label{tab:uav-corridor-raw}
  \centering
  \footnotesize
  \setlength{\tabcolsep}{3pt}
  \begin{tabularx}{\linewidth}{@{}L r rrrrr@{}}
    \toprule
    Model & $\nfe$ & Success & S\&C & \shortstack{Violating\\steps} & Steps & ms/step \\
    \midrule
    \multicolumn{7}{@{}l}{\emph{tilt}}\\
    MeanFM & 1 & $1.000 \pm 0.000$ & $0.000 \pm 0.000$ & $91.3 \pm 34.8$ & $246.4 \pm 7.0$ & $9.4 \pm 0.4$ \\
    MeanFM & 2 & $1.000 \pm 0.000$ & $0.000 \pm 0.000$ & $92.9 \pm 30.2$ & $247.3 \pm 4.8$ & $18.1 \pm 0.4$ \\
    MeanFM & 3 & $1.000 \pm 0.000$ & $0.000 \pm 0.000$ & $92.0 \pm 30.4$ & $246.5 \pm 5.0$ & $27.0 \pm 0.4$ \\
    CI-MeanFM, $\alpha_{\mathrm{end}}=0.2$ & 1 & $1.000 \pm 0.000$ & $0.000 \pm 0.000$ & $90.0 \pm 32.6$ & $242.4 \pm 4.3$ & $9.3 \pm 0.4$ \\
    CI-MeanFM, $\alpha_{\mathrm{end}}=0.2$ & 2 & $1.000 \pm 0.000$ & $0.000 \pm 0.000$ & $93.8 \pm 30.0$ & $246.2 \pm 5.5$ & $18.4 \pm 0.4$ \\
    CI-MeanFM, $\alpha_{\mathrm{end}}=0.2$ & 3 & $1.000 \pm 0.000$ & $0.000 \pm 0.000$ & $93.1 \pm 30.1$ & $245.4 \pm 4.9$ & $27.3 \pm 0.4$ \\
    FM & 1 & $1.000 \pm 0.000$ & $0.000 \pm 0.000$ & $74.5 \pm 25.6$ & $250.9 \pm 5.2$ & $9.1 \pm 0.4$ \\
    FM & 2 & $1.000 \pm 0.000$ & $0.000 \pm 0.000$ & $83.9 \pm 28.4$ & $251.2 \pm 2.9$ & $17.7 \pm 0.4$ \\
    FM & 3 & $1.000 \pm 0.000$ & $0.000 \pm 0.000$ & $87.0 \pm 28.3$ & $251.4 \pm 2.7$ & $25.9 \pm 0.5$ \\
    FM & 5 & $1.000 \pm 0.000$ & $0.000 \pm 0.000$ & $89.9 \pm 28.4$ & $249.7 \pm 2.5$ & $43.4 \pm 0.4$ \\
    FM & 20 & $1.000 \pm 0.000$ & $0.000 \pm 0.000$ & $94.7 \pm 27.0$ & $249.5 \pm 2.9$ & $170.8 \pm 0.4$ \\
    \mainconfig{\baselinemark\ Diffusion} & \mainconfig{20} & \mainconfig{$1.000 \pm 0.000$} & \mainconfig{$0.000 \pm 0.000$} & \mainconfig{$99.4 \pm 27.0$} & \mainconfig{$242.4 \pm 4.2$} & \mainconfig{$178.1 \pm 0.8$} \\
    \midrule
    \multicolumn{7}{@{}l}{\emph{hump}}\\
    MeanFM & 1 & $1.000 \pm 0.000$ & $0.000 \pm 0.000$ & $36.0 \pm 8.9$ & $246.4 \pm 7.0$ & $9.4 \pm 0.4$ \\
    MeanFM & 2 & $1.000 \pm 0.000$ & $0.000 \pm 0.000$ & $35.6 \pm 8.1$ & $247.3 \pm 4.8$ & $18.2 \pm 0.4$ \\
    MeanFM & 3 & $1.000 \pm 0.000$ & $0.000 \pm 0.000$ & $35.7 \pm 7.8$ & $246.5 \pm 5.0$ & $27.2 \pm 0.4$ \\
    CI-MeanFM, $\alpha_{\mathrm{end}}=0.2$ & 1 & $1.000 \pm 0.000$ & $0.000 \pm 0.000$ & $36.6 \pm 8.4$ & $242.4 \pm 4.3$ & $9.3 \pm 0.4$ \\
    CI-MeanFM, $\alpha_{\mathrm{end}}=0.2$ & 2 & $1.000 \pm 0.000$ & $0.000 \pm 0.000$ & $35.6 \pm 7.9$ & $246.2 \pm 5.5$ & $18.0 \pm 0.4$ \\
    CI-MeanFM, $\alpha_{\mathrm{end}}=0.2$ & 3 & $1.000 \pm 0.000$ & $0.000 \pm 0.000$ & $35.6 \pm 7.6$ & $245.4 \pm 4.9$ & $27.1 \pm 0.4$ \\
    FM & 1 & $1.000 \pm 0.000$ & $0.000 \pm 0.000$ & $34.7 \pm 6.8$ & $250.9 \pm 5.2$ & $9.1 \pm 0.4$ \\
    FM & 2 & $1.000 \pm 0.000$ & $0.000 \pm 0.000$ & $32.6 \pm 7.8$ & $251.2 \pm 2.9$ & $17.5 \pm 0.3$ \\
    FM & 3 & $1.000 \pm 0.000$ & $0.000 \pm 0.000$ & $32.6 \pm 8.0$ & $251.4 \pm 2.7$ & $26.3 \pm 0.4$ \\
    FM & 5 & $1.000 \pm 0.000$ & $0.000 \pm 0.000$ & $32.0 \pm 7.9$ & $249.7 \pm 2.5$ & $43.6 \pm 0.4$ \\
    FM & 20 & $1.000 \pm 0.000$ & $0.000 \pm 0.000$ & $31.3 \pm 7.6$ & $249.5 \pm 2.9$ & $172.3 \pm 2.6$ \\
    \mainconfig{\baselinemark\ Diffusion} & \mainconfig{20} & \mainconfig{$1.000 \pm 0.000$} & \mainconfig{$0.000 \pm 0.000$} & \mainconfig{$28.6 \pm 7.7$} & \mainconfig{$242.4 \pm 4.2$} & \mainconfig{$178.2 \pm 2.7$} \\
    \bottomrule
  \end{tabularx}
\end{table}
\dataref{R33 (\texttt{PENDING\_20260923\_uav\_corridor\_v3\_pillars\_v2\_paper\_runs.md}), landed 24-09:
\texttt{temp/23-09-FULL/24-09-1000/batch\_uav\_20260924\_081422/per\_rollout\_detail.csv} and the flown paths
in \texttt{temp/23-09-Corridor-TEMP/plans}, tags \texttt{p23cv3t} (tilt) and \texttt{p23cv3ah} (hump),
CI-MeanFM \texttt{EPlatest\_p23cv3\{t,ah\}}; variant \texttt{diffuser}; the first ten flights of each cell.
v3.83 (author: \enquote{the avoiding in chap6 result is clearly using the 0--1 and the 0.9X $\pm$ range}):
rates over the ten flights with their sample standard deviation, the columns of Table 6.1, success at the
end of the corridor (R45a), steps to it read from the npz on the start-of-step positions. Store:
\texttt{DA\_in\_Paper/data/corridor\_v3\_frontier.json} (\texttt{plotting/extract/corridor\_v3\_frontier.py});
the counts equal \texttt{DA\_20260924\_corridor\_v3.md} \S3.1 and the R45a DA.}

Before projection the cost frontier is drawn over all of them: none satisfies the constraints, so every
configuration is within the band of the best and competes on the two costs of
\autoref{sec:res:avoiding:pareto}, control steps and time per action (\autoref{fig:uav-corridor-tradeoff},
top). Two points are on it. Instantaneous-velocity matching at one evaluation is the cheapest, $250.9$
control steps at $9.1$\,ms, and consistency-interpolated average-velocity matching at one evaluation the
shortest, $242.4$ steps at $9.3$\,ms; it dominates analytic average-velocity matching at the same budget,
which needs $246.4$ steps at $9.4$\,ms. The diffusion baseline flies the corridor in as few steps as the
consistency-interpolated model, at nineteen times its time per action.

\FloatBarrier
\subsubsection{Projection}\label{sec:res:uav:projected}
After projection every flow-based configuration still reaches the end of the corridor on every flight,
except per-step projection at one and two evaluations over the hump, whose flights stop short of the apex
(\autoref{tab:uav-corridor}). Projection buys most of the constraint on this scene, not all of it. The
violating steps of a flight fall from $74.5$ to $99.4$ under the tilt and from
$28.6$ to $36.6$ over the hump to $1.4$ to $10.4$ and $0.2$ to $8.3$, so that in every flow-based row of
the table that reaches the end every flight is free of violations on at least $95\,\%$ of its control
steps; but few are free of them on all. Strictly counted, the only flow-based configuration with violation-free successes is
instantaneous-velocity matching at twenty evaluations under endpoint projection over the hump, $0.800$ of
its flights at $291$\,ms per action. The diffusion baseline, under per-step projection, is violation-free
on most of its flights: $0.800$ under the tilt at $666$\,ms and $1.000$ over the hump at $654$\,ms, each at
its best rule. At the route's goal point, the line these flights were first scored at, the baseline reads
$0.000$ under the tilt and at most $0.600$ over the hump, and every flow-based value is the same
(\autoref{sec:setup:metrics:uav}): $41$ of its $60$ projected flights had passed the end of the corridor and
were still flying towards the goal when the episode limit ended them. The budgets are read as one ladder
across the scene: the per-step projector is active at every budget, endpoint projection at $\nfe\ge 3$,
where it has a guiding step (\autoref{sec:res:constraints:degenerate}), so at $\nfe=1$ and $\nfe=2$ the
table carries per-step rows only.

\begin{table}[htpb]
  \caption[UAV-corridor: the generative models after projection]{UAV-corridor \textbf{after projection}:
  each model at each budget under the per-step projector of \ac{DPCC} and, from $\nfe=3$, under endpoint
  projection, at its best selection rule (most S\&C, then fewest violating steps, then least time; every
  rule in \autoref{app:uav-corridor-rules}); tightened constraints, activation threshold $0.5$, ten flights
  per cell, one training seed. Columns as in \autoref{tab:uav-corridor-raw}; \emph{---}: no flight reaches
  the end of the corridor. At the route's goal point the baseline's S\&C is $0.000$ (tilt) and $0.600$
  under $t$ (hump); every other value is unchanged. \baselinemark\ the diffusion model of \ac{DPCC}.}
  \label{tab:uav-corridor}
  \centering
  \footnotesize
  \setlength{\tabcolsep}{3pt}
  \begin{tabularx}{\linewidth}{@{}L r c rrrrr@{}}
    \toprule
    Model & $\nfe$ & Rule & Success & S\&C & \shortstack{Violating\\steps} & Steps & ms/step \\
    \midrule
    \multicolumn{8}{@{}l}{\emph{tilt, per-step projection}}\\
    MeanFM & 1 & $c$ & $1.000 \pm 0.000$ & $0.000 \pm 0.000$ & $10.2 \pm 1.1$ & $253.2 \pm 4.4$ & $27.4 \pm 0.5$ \\
    MeanFM & 2 & $c$ & $1.000 \pm 0.000$ & $0.000 \pm 0.000$ & $10.3 \pm 0.9$ & $253.9 \pm 4.2$ & $36.4 \pm 0.8$ \\
    MeanFM & 3 & $c$ & $1.000 \pm 0.000$ & $0.000 \pm 0.000$ & $6.6 \pm 0.8$ & $256.4 \pm 2.9$ & $93.6 \pm 1.1$ \\
    CI-MeanFM, $\alpha_{\mathrm{end}}=0.2$ & 1 & $c$ & $1.000 \pm 0.000$ & $0.000 \pm 0.000$ & $10.3 \pm 0.9$ & $252.0 \pm 3.7$ & $27.9 \pm 0.5$ \\
    CI-MeanFM, $\alpha_{\mathrm{end}}=0.2$ & 2 & $c$ & $1.000 \pm 0.000$ & $0.000 \pm 0.000$ & $10.4 \pm 0.7$ & $251.9 \pm 2.4$ & $36.1 \pm 0.3$ \\
    CI-MeanFM, $\alpha_{\mathrm{end}}=0.2$ & 3 & $c$ & $1.000 \pm 0.000$ & $0.000 \pm 0.000$ & $5.0 \pm 0.8$ & $254.0 \pm 1.8$ & $93.4 \pm 1.2$ \\
    FM & 1 & $c$ & $1.000 \pm 0.000$ & $0.000 \pm 0.000$ & $10.2 \pm 1.5$ & $259.8 \pm 9.2$ & $27.6 \pm 1.5$ \\
    FM & 2 & $c$ & $1.000 \pm 0.000$ & $0.000 \pm 0.000$ & $9.8 \pm 1.4$ & $253.9 \pm 9.4$ & $34.7 \pm 0.5$ \\
    FM & 3 & $c$ & $1.000 \pm 0.000$ & $0.000 \pm 0.000$ & $5.8 \pm 1.5$ & $247.2 \pm 3.4$ & $79.2 \pm 3.2$ \\
    FM & 5 & $c$ & $1.000 \pm 0.000$ & $0.000 \pm 0.000$ & $4.4 \pm 1.5$ & $247.6 \pm 2.0$ & $116.9 \pm 2.4$ \\
    FM & 20 & $c$ & $1.000 \pm 0.000$ & $0.000 \pm 0.000$ & $2.1 \pm 0.7$ & $252.6 \pm 2.7$ & $360.4 \pm 1.5$ \\
    \mainconfig{\baselinemark\ Diffusion} & \mainconfig{20} & \mainconfig{$c$} & \mainconfig{$1.000 \pm 0.000$} & \mainconfig{$0.800 \pm 0.422$} & \mainconfig{$0.8 \pm 1.8$} & \mainconfig{$290.3 \pm 15.5$} & \mainconfig{$666.3 \pm 16.7$} \\
    \addlinespace[3pt]
    \multicolumn{8}{@{}l}{\emph{tilt, endpoint projection}}\\
    MeanFM & 3 & $c$ & $1.000 \pm 0.000$ & $0.000 \pm 0.000$ & $7.6 \pm 1.3$ & $255.2 \pm 4.5$ & $73.5 \pm 0.8$ \\
    CI-MeanFM, $\alpha_{\mathrm{end}}=0.2$ & 3 & $c$ & $1.000 \pm 0.000$ & $0.000 \pm 0.000$ & $7.1 \pm 1.0$ & $253.2 \pm 3.5$ & $72.7 \pm 0.5$ \\
    FM & 3 & $c$ & $1.000 \pm 0.000$ & $0.000 \pm 0.000$ & $6.9 \pm 1.8$ & $257.8 \pm 3.3$ & $71.6 \pm 0.6$ \\
    FM & 5 & $c$ & $1.000 \pm 0.000$ & $0.000 \pm 0.000$ & $4.9 \pm 1.9$ & $256.6 \pm 2.9$ & $115.2 \pm 1.4$ \\
    FM & 20 & $c$ & $1.000 \pm 0.000$ & $0.000 \pm 0.000$ & $1.4 \pm 0.7$ & $269.8 \pm 10.5$ & $416.9 \pm 0.8$ \\
    \midrule
    \multicolumn{8}{@{}l}{\emph{hump, per-step projection}}\\
    MeanFM & 1 & $c$ & $0.000 \pm 0.000$ & $0.000 \pm 0.000$ & $0.0 \pm 0.0$ & --- & $26.4 \pm 0.0$ \\
    MeanFM & 2 & $r$ & $0.000 \pm 0.000$ & $0.000 \pm 0.000$ & $0.0 \pm 0.0$ & --- & $35.0 \pm 0.4$ \\
    MeanFM & 3 & $c$ & $1.000 \pm 0.000$ & $0.000 \pm 0.000$ & $4.9 \pm 0.9$ & $286.3 \pm 4.3$ & $91.1 \pm 1.8$ \\
    CI-MeanFM, $\alpha_{\mathrm{end}}=0.2$ & 1 & $t$ & $0.000 \pm 0.000$ & $0.000 \pm 0.000$ & $0.0 \pm 0.0$ & --- & $26.6 \pm 0.3$ \\
    CI-MeanFM, $\alpha_{\mathrm{end}}=0.2$ & 2 & $r$ & $0.000 \pm 0.000$ & $0.000 \pm 0.000$ & $0.0 \pm 0.0$ & --- & $35.4 \pm 0.2$ \\
    CI-MeanFM, $\alpha_{\mathrm{end}}=0.2$ & 3 & $c$ & $1.000 \pm 0.000$ & $0.000 \pm 0.000$ & $4.6 \pm 0.8$ & $282.7 \pm 3.6$ & $89.2 \pm 1.9$ \\
    FM & 1 & $c$ & $0.000 \pm 0.000$ & $0.000 \pm 0.000$ & $0.0 \pm 0.0$ & --- & $26.2 \pm 0.2$ \\
    FM & 2 & $r$ & $0.000 \pm 0.000$ & $0.000 \pm 0.000$ & $0.0 \pm 0.0$ & --- & $34.5 \pm 0.4$ \\
    FM & 3 & $c$ & $1.000 \pm 0.000$ & $0.000 \pm 0.000$ & $5.6 \pm 0.5$ & $277.0 \pm 7.3$ & $77.1 \pm 2.7$ \\
    FM & 5 & $c$ & $1.000 \pm 0.000$ & $0.000 \pm 0.000$ & $3.9 \pm 0.7$ & $277.8 \pm 7.0$ & $114.1 \pm 2.3$ \\
    FM & 20 & $c$ & $1.000 \pm 0.000$ & $0.000 \pm 0.000$ & $1.8 \pm 0.4$ & $285.1 \pm 5.6$ & $355.2 \pm 2.2$ \\
    \mainconfig{\baselinemark\ Diffusion} & \mainconfig{20} & \mainconfig{$r$} & \mainconfig{$1.000 \pm 0.000$} & \mainconfig{$1.000 \pm 0.000$} & \mainconfig{$0.0 \pm 0.0$} & \mainconfig{$320.8 \pm 53.6$} & \mainconfig{$654.2 \pm 89.7$} \\
    \addlinespace[3pt]
    \multicolumn{8}{@{}l}{\emph{hump, endpoint projection}}\\
    MeanFM & 3 & $c$ & $1.000 \pm 0.000$ & $0.000 \pm 0.000$ & $8.0 \pm 1.5$ & $287.7 \pm 4.9$ & $72.1 \pm 0.3$ \\
    CI-MeanFM, $\alpha_{\mathrm{end}}=0.2$ & 3 & $c$ & $1.000 \pm 0.000$ & $0.000 \pm 0.000$ & $8.0 \pm 1.2$ & $284.6 \pm 2.8$ & $71.9 \pm 0.6$ \\
    FM & 3 & $c$ & $1.000 \pm 0.000$ & $0.000 \pm 0.000$ & $8.3 \pm 1.1$ & $299.3 \pm 4.7$ & $71.9 \pm 1.0$ \\
    FM & 5 & $c$ & $1.000 \pm 0.000$ & $0.000 \pm 0.000$ & $3.9 \pm 0.3$ & $292.8 \pm 3.7$ & $112.4 \pm 0.6$ \\
    FM & 20 & single & $1.000 \pm 0.000$ & $0.800 \pm 0.422$ & $0.2 \pm 0.4$ & $316.2 \pm 6.3$ & $291.0 \pm 0.8$ \\
    \bottomrule
  \end{tabularx}
\end{table}
\dataref{R33, the batch, npz and tags of \autoref{tab:uav-corridor-raw}; per-step = \texttt{dpcc-\{r,c,t\}}
at $\eta=0.5$, endpoint = \texttt{hardflow\_sls\{,-r,-c,-t\}}, both with the corridor's stack
\texttt{-bounds\_free-pdes-tightened} (\autoref{sec:setup:protocol:uav}). Success and S\&C at the end of the
corridor, read as the evaluation latches (R45a, \texttt{DA\_20260924\_corridor\_v3\_R45a\_clear\_line.md} \S2;
on recorded start positions only, hump $r$ would read $0.900$ and the best rule be $c$, DA \S3). The best
rule is the same under the near-satisfaction bar for every configuration. Store:
\texttt{data/corridor\_v3\_frontier.json}.}

The budget makes the violations fewer and shallower and does not remove them. Wherever the flights reach
the end of the corridor, the violating steps per flight fall with the budget, under both constraints and
both projection methods. Under the tilt they fall from $9.8$ to $11.1$ at one and two evaluations to $2.1$
to $3.2$ at twenty under per-step projection and to $1.4$ to $2.5$ under endpoint projection, and the
deepest of them from $8.8$ to $11.9$\,cm to $1.3$ to $2.2$ and $1.0$ to $1.3$\,cm. Over the hump they fall
from $4.6$ to $6.5$ under per-step and $8.0$ to $10.4$ under endpoint projection at three evaluations to
$1.8$ to $2.2$ and $0.2$ to $0.7$ at twenty. The budget at which a flight stops reaching the finish line lies
at the other end of the ladder: under per-step projection at one and two evaluations $177$ of the $180$
flow-based flights over the hump stop short of the apex, and the other three reach the end of the corridor
only after drifting across the boundary of its wall. The baseline, at its one budget, is the slowest
planner on the scene: its projected flights need $290$ (tilt) and $321$ (hump) control steps to the end of
the corridor, against $247$ to $270$ and $277$ to $316$ for the flow-based ones, and $623$ to $680$\,ms
per action, the highest price of the scene.

The frontier after projection is drawn at the corridor's near-satisfaction bar
(\autoref{fig:uav-corridor-tradeoff}, bottom): a configuration is in the band when its share of flights that
reach the end of the corridor with at least $95\,\%$ of their control steps free of violations is within
$0.05$ of the best, which every configuration that reaches the end does, and the band then competes on
control steps and time as on D3IL-avoiding. The frontier belongs to the flow-based models, and its cheap end
to the average-velocity objectives. Under the tilt it runs from analytic average-velocity matching at one
evaluation, $253.2$ steps at $27.4$\,ms, through consistency-interpolated average-velocity matching at one
and two evaluations to instantaneous-velocity matching at three evaluations under per-step projection, the
fewest steps, $247.2$, at $79.2$\,ms. Over the hump it is two points: consistency-interpolated
average-velocity matching at three evaluations under endpoint projection, $284.6$ steps at $71.9$\,ms, and
instantaneous-velocity matching at three under per-step projection, $277.0$ at $77.1$\,ms. At the cheapest
budget, then, an average-velocity model dominates instantaneous-velocity matching: under the tilt at one
evaluation it reaches the end in fewer steps at no more time, $253.2$ against $259.8$ at $27.4$ against
$27.6$\,ms, and over the hump at three evaluations with endpoint projection, $284.6$ against $299.3$ at the
same $71.9$\,ms. Under the tilt the margin is smaller than the spread of instantaneous-velocity matching
over its flights, $\pm9.2$ steps; over the hump it is three times the larger of the two spreads. The
baseline is dominated in both panels, with more steps at nine to twenty-four times the time. Two readings
bound this. At the strict count the band holds the baseline alone in both panels: it has the highest share of
violation-free successes, and nothing else comes within $0.05$ of it. At a bar of $97\,\%$ the one- and
two-evaluation cells under the tilt fall out of the band and the tilt's frontier is instantaneous-velocity
matching at three evaluations under per-step projection alone, while over the hump it is unchanged; a bar
below about $87\,\%$ would count unprojected flights over the hump as near-satisfying.

\begin{figure}[htpb]
  \centering
  \includegraphics[width=\linewidth]{fig_uav_corridor_tradeoff}
  \caption[UAV-corridor: steps against cost, before and after projection]{UAV-corridor, control steps to the
  end of the corridor against the time to compute one action, before (a, b) and after (c, d) projection:
  one point per model, budget and projection method at its best rule; colour the model, circles per-step
  projection or none, squares endpoint projection, the number its budget. Rings and the dashed staircase
  mark the non-dominated points of the band, S\&C within $0.05$ of the best before projection, the
  near-satisfaction rate after (\autoref{sec:setup:metrics:uav}); bold labels give strict S\&C above zero.
  (a) and (b) are the same flights. Ten flights per point, one training seed.}
  \label{fig:uav-corridor-tradeoff}
\end{figure}
\dataref{Store figure \texttt{DA\_in\_Paper/figures/da/fig\_uav\_corridor\_tradeoff}
(\texttt{builders/frontier.py::fig\_uav\_corridor\_tradeoff}, v3.83, replacing the archived v2 builder) on
\texttt{data/corridor\_v3\_frontier.json}; frontier rule \texttt{sources.pareto\_front} (band $0.05$), the
avoiding figure's. Near-satisfaction bar and its sensitivity: \texttt{DA\_20260924\_corridor\_v3\_R45a\_clear\_line.md}
Appendix C ($0.90$ gives the same frontier; the floor, unprojected flights $\approx66\,\%$ tilt and $87\,\%$
hump violation-free). At $97\,\%$ (same rule): tilt MeanFM/CI-MeanFM $\nfe=1$ $1/10$, FM $\nfe=1$ $2/10$,
$\nfe=2$ $0$--$2/10$, every $\nfe=3$ per-step cell $10/10$; hump CI-MeanFM and FM $\nfe=3$ endpoint $10/10$,
MeanFM $8/10$. The author asked (24-09) whether MeanFM or CI-MeanFM can be shown to dominate FM: in time
per action alone they tie (DA App.~C); in the frontier of D3IL-avoiding, steps and time at equal
satisfaction, they do at the cheapest budget, as stated. Steps: first start-of-step position past $x=2.0$.
Points coincident on the page are nudged apart horizontally by a few pixels; the data are not.}

Where the violating steps lie is seen in the flights themselves (\autoref{fig:uav-corridor-side}), drawn
for instantaneous-velocity matching, the one flow-based model flown at twenty evaluations, without
projection and under per-step projection at one and at twenty evaluations. Under the tilt the unprojected
flights stay level, so their margin to the leaned plane shrinks steadily along the corridor, to as much
as $0.37$\,m inside the plane at its end. The projected flights move sideways and down until they hold a few
centimetres clear of the plane, and every one of them cuts into it in the last $0.3$\,m of the plane's
span: every violating step of every projected flight, on every model, budget and projection method, lies
between $x=1.71$ and $2.00$\,m. At the first violating step of each projected flow-based flight the
commanded position is already past the end of the plane, at $x=2.27$ to $2.52$\,m and $0.49$ to
$0.61$\,m ahead of the vehicle, and at every violating step the commanded position itself is clear of the
constraint. The plane is enforced over its span of $x$, and on this scene the span is checked at the
commanded position: once that position has left the span, the plan is released from the plane and
returns toward its route, and the vehicle, half a metre behind, follows it into the plane. The violating
steps under the tilt are thus the vehicle catching up with a plan that has already left the constraint,
which is why the budget makes them shallower without removing them. The baseline's projected flights leave
the corridor further to the side of their routes --- a median of $0.37$\,m, against $0.19$\,m for the
flow-based flights --- and stay there, and three to eight of ten of them fly the scene without a
violation.

Over the hump the unprojected flights stay level and pass through the roof. Under per-step projection at one
evaluation every flight climbs along the boundary and stops $0.12$ to $0.13$\,m short of the apex, at
$1.42$ to $1.43$\,m, where it stays until its $396$ control steps run out. At one and two evaluations
all three flow-based models stop this way, on $174$ of their $180$ flights, $0.01$ to $0.17$\,m short of
the apex; the other six clear the apex and then drift sideways across the boundary of the corridor wall.
Two properties of the setup meet at the apex. The plans never rise above $1.30$\,m, the top of the
altitude range of the demonstrations (dotted in \autoref{fig:uav-corridor-side}), while the apex asks for
$1.48$\,m. And each face of the roof, like the tilt, is enforced over its span of $x$ checked at the
commanded position, so while that position is short of the apex the rising face binds the whole plan,
including its part beyond the apex, where the face keeps rising. Both hold at every budget. At one and
two evaluations the corrected plan settles short of the apex; from three on it does not, and at twenty
the flights cross the apex grazing the boundary by about a centimetre and descend toward their route
altitude before the line.

\begin{figure}[htpb]
  \centering
  \includegraphics[width=\linewidth]{fig_uav_corridor_side}
  \caption[UAV-corridor: the flights along the corridor]{UAV-corridor along the corridor,
  instantaneous-velocity matching, ten flights per series, one training seed: without projection at
  $\nfe=1$ (grey) and under per-step projection with cumulative-cost selection at $\nfe=1$ and $\nfe=20$
  (blue). (a, b) Tilt: the margin of the flown position to the leaned plane, less the rotor reach, over
  the plane's span of $x$; a control step below the dashed line violates. (c, d) Hump: altitude over $x$,
  the roof shaded, the violation boundary dashed, the highest planned altitude dotted; a circle ends a
  flight short of the finish line.}
  \label{fig:uav-corridor-side}
\end{figure}
\dataref{Store figure \texttt{DA\_in\_Paper/figures/da/fig\_uav\_corridor\_side} (builder
\texttt{builders/altitude.py::fig\_uav\_corridor\_side}; extract \texttt{extract/corridor\_v3\_side.py}
$\to$ \texttt{data/corridor\_v3\_side.json}), cells read through \texttt{analysis/corridor\_v3\_grid.py}:
FM, \texttt{diffuser} at $\nfe=1$ and \texttt{dpcc-c-bounds\_free-pdes-tightened} at 1 and 20, the first
ten flights; $c$ is FM's best per-step rule at both budgets and both constraints (\autoref{tab:uav-corridor}).
Margin $=\mathbf{n}\cdot(\mathbf{p}-\mathbf{p}_0)-0.31$ with the evaluation's own \texttt{\_hs\_normal3};
hump boundary $=$ roof $+\,0.31\sqrt{1+s^2}$, $s=1.10/1.5$; highest planned altitude $1.2992$\,m, the
largest planned $z$ of the drawn cells. Violating-step positions of all $560$ projected tilt flights,
$4115$ steps: $1.709\le x\le 2.000$, $99\,\%$ at $x\ge 1.75$ (the DA's $[1.75,2.00]$ is its 5--95\,\%
range). Commanded position at the first violating step, $530$ projected flow-based flights: $x=2.27$--$2.52$
(median $2.38$), $0.49$--$0.61$\,m ahead (median $0.55$); clear of the constraint at $4014$ of $4014$
violating steps. Span check at the commanded position under \texttt{-pdes}:
\texttt{eval\_mix\_uav.py:1830--1833} (\texttt{\_gate\_x = p\_des[0]}) and \texttt{:1483--1484}. Hump
stops: per flight $x=-0.174$ to $-0.006$, $z=1.38$ to $1.51$\,m, every flight on the $396$-step cap; the six
others are MeanFM and CI-MeanFM at $\nfe=2$ under $t$, three each, wall violations with no collision
(\texttt{phys\_contact\_frac} $0$), ending $0.95$--$1.79$\,m from the goal. Plan bound: every stored plan
$z\in[0.9008,1.2992]$\,m (DA \S4.2, \texttt{LimitsNormalizer.unnormalize}); a third cause the DA names ---
the correction is the smallest in normalised coordinates --- is not stated in the text. Replaces the
three placeholders \texttt{fig:uav-corridor-tradeoff}, \texttt{-paths}, \texttt{-altitude} (v3.80).}

\paragraph{Altitude.} Both constraints of this scene act on altitude, and that is what they are for: under
the tilt the room is below the launch altitude, so a projected flight has to descend as it moves sideways;
under the hump it has to climb. Unprojected, every model flies level, between $0.83$ and $1.27$\,m where
the constraints act. Projected, the lowest altitude under the tilt drops on $540$ of the $560$ projected
flights, by a median of $0.12$ to $0.37$\,m per configuration, and the highest altitude over the hump
rises on all $560$, by a median of $0.27$ to $0.39$\,m and on single flights by up to $0.60$\,m
(\autoref{tab:uav-corridor-altitude}). Under the tilt the descent deepens with the budget, from $0.12$ to
$0.19$\,m for the two average-velocity models at one and two evaluations to $0.21$ to $0.26$\,m at three,
and from $0.19$ to $0.25$\,m for instantaneous-velocity matching at one and two to $0.35$ to $0.37$\,m at
twenty; the climb over the hump does not.

\begin{table}[htpb]
  \caption[UAV-corridor: the altitude change under projection]{UAV-corridor, the altitude change that
  projection makes: under the tilt the lowest altitude over $0.5\le x\le 2.0$\,m, over the hump the
  highest over $-0.5\le x\le 0.5$\,m, each the projected flight minus the unprojected flight of the same
  trial, as the median over a cell's ten flights. The range runs over the model's projected cells: every
  budget, projection method and selection rule. One training seed. \baselinemark\ the diffusion model of
  \ac{DPCC}.}
  \label{tab:uav-corridor-altitude}
  \centering
  \footnotesize
  \setlength{\tabcolsep}{5pt}
  \begin{tabularx}{\linewidth}{@{}L l rr@{}}
    \toprule
    Model & $\nfe$ & \shortstack{tilt: lowest altitude,\\median change (m)} & \shortstack{hump: highest altitude,\\median change (m)} \\
    \midrule
    MeanFM & 1, 2, 3 & $-0.26$ to $-0.13$ & $+0.33$ to $+0.39$ \\
    CI-MeanFM, $\alpha_{\mathrm{end}}=0.2$ & 1, 2, 3 & $-0.26$ to $-0.12$ & $+0.33$ to $+0.39$ \\
    FM & 1, 2, 3, 5, 20 & $-0.37$ to $-0.19$ & $+0.27$ to $+0.36$ \\
    \mainconfig{\baselinemark\ Diffusion} & \mainconfig{20} & \mainconfig{$-0.30$ to $-0.22$} & \mainconfig{$+0.32$ to $+0.34$} \\
    \bottomrule
  \end{tabularx}
\end{table}
\dataref{DA \S3.6 and Appendix B (median per cell and rule), recomputed from the npz at v3.80 with
\texttt{corridor\_v3\_grid.altitude}; flights paired by trial index. $540$ of $560$ tilt flights descend
(the $20$ that do not: MeanFM $8$, CI-MeanFM $10$, FM $2$, each within $+0.11$\,m); $560$ of $560$ hump
flights climb. The unprojected band is the range of the same window extremes. The v3.65 pilots
($-0.15$ to $-0.29$, $+0.30$ to $+0.40$) lie inside these ranges.}

UAV-pillars is not drawn as a cost frontier: its result is the difference between a plan's outcome on
D3IL-avoiding and the outcome of the same plan flown, which is a comparison and not a trade-off
(\autoref{sec:res:uav:models}). UAV-s-curve has data, but no configuration flown on it succeeds with
constraint satisfaction (\autoref{tab:uav-scurve}); a cost frontier over points that all share one outcome
is a ranking of failures by price, so that scene gets the abort count printed beside each of its rows
instead.

\FloatBarrier
\subsubsection{Projection Methods}\label{sec:res:uav:projection}
Per-step and endpoint projection separate on this scene only at twenty evaluations
(\autoref{tab:uav-corridor}; every selection rule in \autoref{app:uav-corridor-rules}). At three and five evaluations the two are level, at no
violation-free success under either constraint; endpoint projection is as cheap or cheaper, $44$ to
$74$\,ms per action against $77$ to $94$ at three evaluations and $73$ to $115$ against $112$ to $118$ at
five, and it leaves more violating steps: at three evaluations about two more per flight under the tilt
($6.9$ to $9.2$ against $5.0$ to $7.5$) and about four more over the hump ($8.0$ to $10.4$ against $4.6$
to $6.5$), at five about one more. At twenty evaluations endpoint projection leaves fewer and shallower
violating steps than per-step projection under the tilt, $1.4$ to $2.5$ per flight and at most
$1.3$\,cm deep against $2.1$ to $3.2$ and $2.2$\,cm, and
over the hump it is the only flow-based configuration that succeeds violation-free, on four to eight
flights of ten depending on the rule, where per-step projection succeeds on none. The bar this thesis set
for endpoint projection, to overtake the per-step projector of \ac{DPCC} at a lower cost, is met on one
constraint at one budget: over the hump at twenty evaluations with a single candidate, eight
violation-free successes against none, at $291$ against $355$\,ms per action. With four candidates it
still overtakes the per-step projector there, at a higher cost ($422$ to $425$\,ms).

\guard{Endpoint projection has no row at $\nfe=1$ or $\nfe=2$ on this scene and will not get one: at the
activation threshold of $0.5$ the only active sampling step there is the terminal one, so there is no
guiding step for the method to act on (\autoref{sec:res:constraints:degenerate}). That is a property of
the method at those budgets, not a pending evaluation. The baseline has no endpoint rows for another
reason: endpoint projection adds its control to a velocity field (\autoref{sec:method:hardflow}), which
the diffusion model does not have.}

The projection configuration of this scene differs from that of the other quadrotor scene
(\autoref{sec:setup:protocol:uav}), so the scenes are not pooled.

\FloatBarrier
